\documentclass[11pt]{article}
\usepackage{icmlko}
\begin{document}

\title{원핫의 0 은 무지의 정직한 표시였다 --- 모르는 걸 계산해 주자
더 나빠졌다}
\author{월드모델 실험실}
\date{노트 442 · 2026년 8월 2일}
\maketitle

\begin{abstract}
노트 441 이 채택한 것은 \emph{휴전}이었다 --- 안 배운 수집에는
\texttt{entry\_friction} 을 \textbf{안 쓴다}. 그런데 그 축은 도메인
\textbf{안}에서는 값이 크다(모바일 $+0.648$ · 도서 $+0.273$). 못 쓰는
이유는 정보가 없어서가 아니라 \textbf{어느 쪽으로 읽을지를 모르기}
때문이고, 집 안에서 그걸 알려 주는 원핫과 grp 이 \textbf{집 밖에는 둘 다
없다}(원핫은 전부 $0$, grp 은 값 자체가 없다).

그래서 \textbf{새 수집이 자기 행으로 계산할 수 있는 표지}를 줘 봤다 ---
그 수집의 \texttt{entry\_friction} 평균과 SD. 새 수집도 자기 행에서 바로
계산되니 원핫과 달리 \textbf{집 밖에서도 값이 있다}. 등록한 문장은
``가리지 않고도 가림보다 낫거나 같다'' 였다.

\textbf{반증됐다. 그것도 크게.} 표지를 주자 두 짝 다 \emph{더} 나빠졌다.

\begin{center}
\begin{tabular}{lrr}
\hline
팔 & KR 만화 & 비게임 앱 \\
\hline
가 · 기준 (가림 없음) & $+0.3206$ & $-0.0501$ \\
나 · 가림 (노트 441 이 채택) & $\mathbf{+0.5336}$ & $\mathbf{+0.2635}$ \\
\textbf{다 · 평균표지} & $+0.2033$ & $\mathbf{-0.2905}$ \\
라 · 평균표지 $+$ 가림 & $+0.5450$ & $+0.2802$ \\
\hline
\end{tabular}
\end{center}

\textbf{앱은 $-0.050$ 에서 $-0.291$ 로 갔다.} 아무것도 안 준 것보다
\emph{나쁘다}.
\end{abstract}

\begin{figure}[t]
\centering
\includegraphics[width=.86\textwidth]{figs/f1.pdf}
\caption{두 집 밖 짝에서 네 팔. 초록(가림)이 노트 441 이 채택한 것이고,
빨강(평균표지)이 이번에 시험한 대안이다. 대안은 기준보다도 아래로
내려간다. 파랑(표지 $+$ 가림)이 초록과 거의 같은 것이 진단이다 ---
\textbf{표지가 해로운 것은 오직 독 있는 축이 켜져 있을 때}다.}
\end{figure}

\section{왜 나빠지나 --- 마지막 줄이 진단이다}

네 번째 팔이 답을 준다. \textbf{가림을 같이 걸면 평균표지는 무해해진다}
($+0.5450$ · $+0.2802$ --- 가림만 한 것과 사실상 같다). 즉 표지 자체가
독인 게 아니라 \textbf{표지가 독 있는 축을 쓰게 만든다}.

기계는 이렇다. 상수 열은 노트 412 가 말한 대로 \textbf{도메인 이름표를 한
번 더 적는 것}이다. 나무는 그 값으로 갈라 \emph{도메인마다 다른 규칙}을
학습한다. 집 안에서는 원핫과 겹쳐 아무것도 안 늘지만, \textbf{집 밖에서는
겹칠 원핫이 없다}. 새 수집의 평균값은 학습한 어느 도메인 근처에
떨어지고, 모델은 \textbf{그 도메인의 규칙을 자신 있게 갖다 쓴다}. KR 의
평균 $0.268$ 은 만화($0.112$)와 펀딩($0.396$) 사이에, 앱의 $0.295$ 도
같은 자리에 떨어진다 --- 그리고 그 자리의 규칙은 KR 의 부호($-0.320$)와
반대다.

\textbf{원핫이 전부 $0$ 인 것은 결함이 아니라 성질이었다.} 그것은
``모른다''를 정직하게 적는다. 계산되는 표지는 ``이 근처다''라고 적고,
근처가 틀리면 모델은 \emph{틀린 쪽을 확신}한다. 노트 399 이래 원핫을
여러 번 범인으로 몰았는데, 집 밖에서는 원핫의 $0$ 이 오히려 안전장치다.

\section{채택 판정}

\textbf{아무것도 안 바꾼다.} 팔 라(표지 $+$ 가림)가 팔 나보다 두 짝에서
$+0.0114$ · $+0.0167$ 높지만, 이 크기는 집 밖 짝 둘의 차이를 가를 자가
없다 --- 노트 441 이 만든 규약대로 집 밖 효과는 집 밖 짝으로 재는데,
짝이 둘뿐이고 차이가 채택된 효과($+0.21 \sim +0.31$)의 5\% 다. 판은 이
변경도 못 본다. \textbf{유지}: ``예보 때 \texttt{entry\_friction} 만
가리는 F18''.

\section*{남은 것}

\begin{itemize}
\item \textbf{어느 쪽으로 읽을지를 알려 주는 다른 방법.} 계산되는 표지는
죽었다. 남은 후보는 \textbf{부호 자체를 새 수집에서 추정}하는 것 ---
라벨 몇 개만 있으면 짝 안 부호를 재서 축을 뒤집을 수 있다. 그건
\emph{완전한} 영샷이 아니라 \textbf{몇 샷}이다. 값이 얼마나 되는지,
라벨 몇 개면 되는지는 안 재 봤다.
\item \textbf{CN 이 왜 안 오르나} --- 독을 안 먹었는데 $+0.2249$ 다.
라벨 분포(SD $0.485$ · 왜도 $1.84$)가 마지막 후보다.
\item LODO 씨앗 셋$\to$아홉 --- 다섯 노트째 미룸.
\end{itemize}

\end{document}
